\section{System Context and Methodology}\label{sec:method}

\subsection{The system under study}\label{sec:system}

The subject system (public repository \sysname{}) is a two-layer
middleware that connects self-hosted and commercial LLMs to OpenClaw, an
open-source personal-agent framework bridging WhatsApp and Discord. It
has been in continuous production on a single macOS host since early
March 2026. Architecturally it follows a three-plane design
(Fig.~\ref{fig:planes}):

\begin{itemize}
\item \textbf{Control plane} --- a tool-governance proxy (tool filtering
to a hard cap, schema repair, alert-context stripping, custom tool
interception), a declarative governance engine (90~invariants,
827~checks, 23~meta-rules, 15~mechanized discovery scanners, daily audit
cron), SLO monitoring, circuit breakers, and a convergence engine that
machine-synchronizes declared state (job registry, provider registry,
service registry) to runtime state (crontab, fallback chains, launchd).
\item \textbf{Capability plane} --- an adapter routing chat/completions
across 8 providers (self-hosted Qwen3-235B primary, Qwen2.5-VL multimodal
route, plus commercial fallbacks) with capability-scored automatic
fallback chains.
\item \textbf{Memory plane} --- a knowledge base ($\sim$1{,}100+ notes,
local-embedding RAG index), multimodal media index, conversation-harvest
pipeline, and daily LLM-driven synthesis jobs (``dream'', ``evening
digest'', ``deep dive'') that push insights to the user.
\end{itemize}

\begin{figure}[t]
\centering
\begin{tikzpicture}[node distance=4mm and 5mm]
  % user layer
  \node[box, fill=blue!6, minimum width=118mm] (user)
    {User layer: WhatsApp + Discord (push output, real human reader)};
  % planes
  \node[stage, minimum width=38mm, minimum height=26mm, below=7mm of user.south west, anchor=north west, align=center]
    (ctrl) {\textbf{Control plane}\\[1pt]\scriptsize tool governance proxy\\\scriptsize 90 invariants / 827 checks\\\scriptsize convergence engine\\\scriptsize SLO + circuit breakers};
  \node[stage, minimum width=38mm, minimum height=26mm, right=4mm of ctrl]
    (cap) {\textbf{Capability plane}\\[1pt]\scriptsize 8-provider adapter\\\scriptsize multimodal routing\\\scriptsize capability-scored\\\scriptsize fallback chains};
  \node[stage, minimum width=38mm, minimum height=26mm, right=4mm of cap]
    (mem) {\textbf{Memory plane}\\[1pt]\scriptsize KB + RAG index\\\scriptsize conversation harvest\\\scriptsize daily LLM synthesis\\\scriptsize (dream / digest / deep dive)};
  % observation band
  \node[con, minimum width=118mm, below=7mm of cap.south, anchor=north]
    (obs) {\scriptsize Observation band: 4{,}286 unit tests ~$\cdot$~ daily governance audit ~$\cdot$~ preflight ~$\cdot$~ watchdogs ~$\cdot$~
           \textbf{weekly human user-view ritual} ~$\cdot$~ LLM observer};
  \draw[flow] (user.south -| ctrl.north) -- (ctrl.north);
  \draw[flow] (user.south -| cap.north) -- (cap.north);
  \draw[flow] (user.south -| mem.north) -- (mem.north);
  \draw[flow, dashed] (obs.north -| ctrl.south) -- (ctrl.south);
  \draw[flow, dashed] (obs.north -| cap.south) -- (cap.south);
  \draw[flow, dashed] (obs.north -| mem.south) -- (mem.south);
\end{tikzpicture}
\caption{The system under study: a three-plane agent runtime with an
observation band beneath it. The incidents in this paper are, by
selection, the ones that crossed all planes \emph{and} the observation
band without raising an actionable signal---most were finally caught at
the very top, by the human reading pushed output
(\S\ref{sec:discovery}).}
\label{fig:planes}
\end{figure}

Scale snapshot at the study cutoff (2026-06-11): ${\sim}40$ registered
scheduled jobs across two cron schedulers; 8 LLM providers; 3 supervised
long-running services; 4{,}286 unit tests in 121 suites; 22 published
incident postmortems. One human operator (the system owner) and one AI
engineering collaborator (Claude, used through a coding-agent interface)
develop and operate the system; the agent runtime itself is powered by
separate models (Qwen3-class), making this---to our knowledge---also a
data point on AI-assisted operation of an AI system.

\subsection{Incident corpus and postmortem protocol}\label{sec:protocol}

The corpus is every incident between 2026-04-09 and 2026-06-02 that
(a)~reached production, (b)~had a silent phase---a period in which the
system was failing or had failed while all automated indicators stayed
green---and (c)~was closed with a full postmortem. 22 incidents qualify.
Postmortems follow a mandatory in-repo protocol (the
``exception-analysis constitution'') requiring, before any fix:

\begin{enumerate}
\item \textbf{Full causal-chain diagram} --- timeline $\times$ layer
$\times$ logic $\times$ architecture, from upstream trigger to
user-perceived symptom, with code-branch truth values annotated;
\item \textbf{Three-layer root cause} --- \emph{trigger} (what external
event ignited it), \emph{amplifier} (what architectural flaw spread it),
\emph{concealer} (what absence hid it until a human noticed);
\item \textbf{Timeline reconstruction} with per-minute precision where
logs allow;
\item \textbf{Condition-combination analysis} --- why now and not before:
which set of individually benign, often long-latent conditions first
co-occurred;
\item \textbf{Feeding the governance ontology} --- new invariants,
meta-rule candidacy, catalog entry.
\end{enumerate}

The protocol was itself a product of early incidents (it was adopted
after the second incident in the corpus) and was applied retroactively to
the first two. Each postmortem is a standalone document in the public
repository, and the consolidated catalog is the canonical index from
which this paper's taxonomy is built.

\subsection{Taxonomy construction}\label{sec:construction}

We classify by \textbf{mechanism} (how the error evaded observation)
rather than by \textbf{location} (which job or file failed). The
location-oriented draft of the catalog was produced first and discarded:
the same mechanism (e.g., positional parsing of LLM output; copy-pasted
error-suppression idioms) recurred across unrelated jobs, so location
classes had no predictive or defensive value, while one mechanism-level
defense (e.g., a repo-wide scanner for the suppression idiom) immunized
every location at once. Classes were derived iteratively: each new
postmortem was matched against existing classes' ``true root cause''
column; two consecutive non-matches forced a class split or a new class.
The five-class scheme has been stable since mid-May 2026 over the last 9
incidents.

Two limitations are flagged here and expanded in
\S\ref{sec:threats}: classification was performed by the two system
operators (human + AI collaborator) without independent annotators, so we
report no inter-annotator agreement; and selection requires a
\emph{completed} postmortem, so failures still silent today are by
construction absent (we discuss this surviving-silence bias in
\S\ref{sec:threats}).

\subsection{The meta-pattern counter}\label{sec:counter}

Early in the study we adopted a meta-rule, \textbf{MR-4
``silent-failure-is-a-bug''}, and began counting its
\emph{manifestations}---distinct incidents or sub-events in which an
error signal existed somewhere in the system but never reached a human in
actionable form. The counter reached ${\geq}28$ across the 22 incidents
(several incidents contain multiple distinct manifestations; numbering
has deliberate gaps where manifestations were recorded inside hotfix
notes). We report the counter not for its precision but for its shape:
manifestations continued to appear at a roughly constant rate even as
defenses accumulated---in \emph{new forms} each time
(\S\ref{sec:bugclass}). This is the empirical basis for treating silent
failure as a bug class rather than a finite list of bugs.
